\documentclass[14pt,a4paper]{extarticle}

\usepackage[left=20mm, right=15mm, top=15mm, bottom=15mm]{geometry}
\usepackage{polyglossia}
\setdefaultlanguage{russian}
\setotherlanguage{english}
\defaultfontfeatures{Mapping=tex-text}
\usepackage{graphicx}
\usepackage{subcaption}
\usepackage{amsmath}
\usepackage{amssymb}
\usepackage{hyperref}
\usepackage{listings}
\usepackage{xcolor}
\usepackage{caption}
\usepackage{setspace}
\usepackage{indentfirst}
\usepackage{titlesec}
\usepackage{tocloft}
\usepackage{longtable}
\usepackage{booktabs}
\usepackage{fancyhdr}
\usepackage{tcolorbox}
\captionsetup[figure]{position=below, labelsep=period, font={small,it}}
\captionsetup[table]{position=above, labelsep=period, font={small,it}}

\usepackage{fontspec}
\setmainfont{Linux Libertine}
\newfontfamily\cyrillicfont{Linux Libertine}[Script=Cyrillic]
\newfontfamily\cyrillicfontsf{Linux Libertine}[Script=Cyrillic]
\newfontfamily\cyrillicfonttt{DejaVu Sans Mono}[Script=Cyrillic]
\setmonofont{DejaVu Sans Mono}
\setstretch{1.15}
\titleformat{\section}{\normalfont\bfseries\fontsize{16}{24}\selectfont}{\thesection}{1em}{}
\titleformat{\subsection}{\normalfont\bfseries\fontsize{16}{24}\selectfont}{\thesubsection}{1em}{}

\lstset{
  basicstyle=\small\ttfamily,
  breaklines=true,
  frame=single,
  captionpos=t,
  numbers=left,
  numberstyle=\tiny,
  tabsize=4,
  showstringspaces=false,
  keywordstyle=\color{blue},
  commentstyle=\color{gray},
  stringstyle=\color{orange},
}

\hypersetup{colorlinks=true, linkcolor=black, urlcolor=blue, citecolor=black}
\pagestyle{fancy}
\fancyhf{}
\fancyfoot[C]{\thepage}
\renewcommand{\headrulewidth}{0pt}
\setlength{\parindent}{1.25cm}


\begin{document}

\begin{titlepage}

\begin{center}
	{\large
МИНИСТЕРСТВО НАУКИ И ВЫСШЕГО ОБРАЗОВАНИЯ РФ

		РГУ им. А.Н. Косыгина \\ (Технологии. Дизайн. Искусство)

		\bigskip

		Институт Информационных технологий и цифровой трансформации
	}

	\bigskip



	\textbf{Кафедра <<Искусственного интеллекта, прикладной математики и программирования>>}
	\bigskip
	\vfill \textsc{\Large Курсовая работа} \\
	{\large по дисциплине <<Вероятностное моделирование процессов и систем>>}

	\bigskip

	на тему: Создание приложений вероятностного моделирования
\end{center}

\vspace*{1.5cm}

\hfill
\begin{minipage}{.6\linewidth}
	\begin{tabular}{c}
		\textbf{Выполнил:} студент группы ИТПМ-124    \\\hline \\[.3cm]
		{\large Фриев Д.С.}                       \\ \hline \scriptsize{(Фамилия, имя, отчество)}
		\\[.3cm] \\ \hline
		\scriptsize{(подпись)}                       \\[.3cm]
		\textbf{Принял: }\hfill доцент кафедры ИИПМиП \\\hline
		\\[.3cm]
		Романенков А.М.                             \\ \hline
		\scriptsize{(Фамилия, имя, отчество)}
		\\[.3cm] \\ \hline
		\scriptsize{(подпись)}
	\end{tabular}
	\vspace*{1cm}
\end{minipage}

\centerline{\textbf{Оценка}: \hspace*{8cm} \textbf{Дата}: \hspace*{2cm}}
\vspace*{1cm}
\centerline{Москва, 2026}

\end{titlepage}
\clearpage
\tableofcontents
\clearpage

\section*{Введение}
\addcontentsline{toc}{section}{Введение}

\subsection{Цель и задачи}

Цель работы — программная реализация комплекса задач вероятностного моделирования. Для достижения цели необходимо решить следующие задачи:


\begin{itemize}
  \item разработать математические модели для каждого задания;
  \item продумать и создать архитектуру приложений;
  \item реализовать пользовательский интерфейс для управления параметрами моделирования;
\end{itemize}

\subsection{Используемые технологии}

Программная реализация выполнена на языке \texttt{C++} стандарта \texttt{C++23} с использованием следующих библиотек и инструментов:


\begin{itemize}
  \item \texttt{CMake} (\url{https://github.com/Kitware/CMake}) — система сборки проекта;
  \item \texttt{Qt6 Quick} (\url{https://github.com/qt/qtdeclarative}) — фреймворк для разработки графического интерфейса;
  \item \texttt{GLFW} (\url{https://github.com/glfw/glfw}), \texttt{Dear ImGui} (\url{https://github.com/ocornut/imgui}) и \texttt{ImPlot} (\url{https://github.com/epezent/implot}) — для визуализации графиков дискретной случайной величины в одном из заданий;
  \item \texttt{nlohmann::json} (\url{https://github.com/nlohmann/json}) — парсинг конфигурационных файлов;
  \item \texttt{kursach-autogen} (\url{https://github.com/koftamainee/kursach-autogen/}) — генерация пояснительной записки.
\end{itemize}

Все разработанные приложения устойчивы к аварийным завершениям, предусмотрена обработка ошибок ввода-вывода и некорректных параметров.


\subsection{Структура пояснительной записки}

Документ имеет следующую структуру:

\begin{itemize}
  \item введение;
  \item описание задач с подробностями их реализации;
  \item выводы;
  \item список использованных источников;
  \item приложение со ссылкой на репозиторий исходного кода.
\end{itemize}

\clearpage

\section{Выбор символа из алфавита}

Из упорядоченного алфавита, содержащего $N$ различных символов, последовательно
выбирают два символа: сначала один, затем — один из оставшихся. Все исходы считаются
равновероятными. Требуется определить вероятность того, что в первый раз, во второй раз
и оба раза будет выбран символ с чётным номером. Индексация элементов алфавита начинается
с $1$. Рассматриваются два варианта нумерации после первой выборки:
а) индексы оставшихся элементов не изменяются;
б) индексы изменяются (переиндексируются начиная с $1$).
Необходимо реализовать приложение для проведения эксперимента, которое принимает через
командную строку мощность алфавита $N$ и количество экспериментов $C$, и выполнить
программное сравнение эмпирической и аналитической вероятностей.


\subsection{Аналитическое решение}

Пусть алфавит содержит $N$ символов с номерами от $1$ до $N$. Число чётных номеров равно
$\lfloor N/2 \rfloor$. Первый символ выбирается равновероятно из всех $N$ элементов.


\textbf{Случай а) — фиксированные индексы.}

\begin{equation}\label{eq.p_first_fixed}
P(\text{первый чётный}) = \frac{\lfloor N/2 \rfloor}{N}
\end{equation}

При фиксированных индексах после удаления первого элемента оставшиеся $N-1$ символов
сохраняют свои номера. Вероятность того, что второй выбранный символ имеет чётный номер,
зависит от чётности первого:


\begin{equation}\label{eq.p_second_fixed}
P(\text{второй чётный}) = \frac{\lfloor N/2 \rfloor}{N} \cdot \frac{\lfloor N/2 \rfloor - 1}{N-1} + \frac{\lceil N/2 \rceil}{N} \cdot \frac{\lfloor N/2 \rfloor}{N-1}
\end{equation}

\begin{equation}\label{eq.p_both_fixed}
P(\text{оба чётных}) = \frac{\lfloor N/2 \rfloor}{N} \cdot \frac{\lfloor N/2 \rfloor - 1}{N - 1}
\end{equation}

\textbf{Случай б) — переиндексация после первой выборки.}

После удаления первого элемента оставшиеся $N-1$ символов получают новые порядковые
номера от $1$ до $N-1$. Число чётных позиций среди $N-1$ элементов равно $\lfloor (N-1)/2 \rfloor$.
Вероятность чётного номера при первом выборе не изменяется:


\begin{equation}\label{eq.p_first_reindexed}
P(\text{первый чётный}) = \frac{\lfloor N/2 \rfloor}{N}
\end{equation}

После переиндексации вероятность чётного номера при втором выборе не зависит от результата
первого, поскольку новая нумерация определяется позицией в оставшемся списке:


\begin{equation}\label{eq.p_second_reindexed}
P(\text{второй чётный}) = \frac{\lfloor (N-1)/2 \rfloor}{N-1}
\end{equation}

\begin{equation}\label{eq.p_both_reindexed}
P(\text{оба чётных}) = \frac{\lfloor N/2 \rfloor}{N} \cdot \frac{\lfloor (N-1)/2 \rfloor}{N-1}
\end{equation}

\subsection{Эмпирическая проверка}

Программа принимает два параметра командной строки: мощность алфавита $N$ и количество
экспериментов $C$. Каждый эксперимент моделирует два последовательных выбора символа.
Для генерации случайных чисел используется Вихрь Мерсенна (\texttt{std::mt19937}).
Запуск серии экспериментов выполняется через вспомогательный класс \texttt{TaskRunner},
который принимает произвольный эксперимент в виде callable и возвращает вектор результатов.


\begin{lstlisting}[language=C++, caption={Эксперимент с фиксированными индексами}, label=lst.task01_fixed]
auto experiment_fixed_indices = [alphabet_size](auto& rng) {
    std::uniform_int_distribution<size_t> dist1(1, alphabet_size);
    size_t first = dist1(rng);

    std::vector<size_t> remaining;
    remaining.reserve(alphabet_size - 1);
    for (size_t i = 1; i <= alphabet_size; i++) {
        if (i != first) remaining.push_back(i);
    }

    std::uniform_int_distribution<size_t> dist2(0, remaining.size() - 1);
    size_t second = remaining[dist2(rng)];

    return Outcome{first % 2 == 0, second % 2 == 0};
};
\end{lstlisting}

\begin{lstlisting}[language=C++, caption={Эксперимент с переиндексацией после первого выбора}, label=lst.task01_reindexed]
auto experiment_reindexed = [alphabet_size](auto& rng) {
    std::vector<size_t> indices(alphabet_size);
    std::iota(indices.begin(), indices.end(), 1);

    std::uniform_int_distribution<size_t> dist1(0, indices.size() - 1);
    size_t first_pos = dist1(rng);
    bool first_even = (first_pos + 1) % 2 == 0;

    indices.erase(indices.begin() + first_pos);

    std::uniform_int_distribution<size_t> dist2(0, indices.size() - 1);
    size_t second_pos = dist2(rng);
    bool second_even = (second_pos + 1) % 2 == 0;

    return Outcome{first_even, second_even};
};
\end{lstlisting}

Для подтверждения корректности полученных аналитических выражений был проведён
вычислительный эксперимент при фиксированном размере алфавита $N = 100$ и количестве
экспериментов $C = 10^8$ (сто миллионов итераций). Такой объём выборки обеспечивает
высокую точность оценки вероятностей и позволяет надёжно проверить сходимость
эмпирических частот к теоретическим значениям.


Результаты эксперимента представлены на рис.~\ref{fig.convergence}.


\begin{figure}[h]
\centering
\includegraphics[width=0.9\linewidth]{assets/task01.png}
\caption{Результаты эксперимента при $N = 100$, $C = 10^8$}\label{fig.convergence}
\end{figure}

Полученные эмпирические вероятности:


\begin{table}[h]
\centering
\caption{Результаты эксперимента при $N = 100$, $C = 10^8$}\label{tbl.results_task01}
\begin{tabular}{l | l | l | l}
\hline
Вариант нумерации & Первый чётный & Второй чётный & Оба чётных \\
\hline
Фиксированные индексы & 0.499913 & 0.500032 & 0.247482 \\
Переиндексация & 0.499995 & 0.494859 & 0.247438 \\
\hline
\end{tabular}
\end{table}

Теоретические значения при $N = 100$ составляют:


\[
P_{\text{first}} = \frac{50}{100} = 0.5
\]

\[
P_{\text{second}}^{(a)} = 0.5
\]

\[
P_{\text{second}}^{(b)} = \frac{49}{99} \approx 0.494949
\]

\[
P_{\text{both}} = \frac{50}{100} \cdot \frac{49}{99} \approx 0.247475
\]

Сравнение эмпирических и теоретических значений показывает высокую точность
моделирования. Максимальное абсолютное отклонение составляет $8.7 \times 10^{-5}$
для случая фиксированных индексов и $9.0 \times 10^{-5}$ для случая переиндексации.
Таким образом, экспериментальные данные подтверждают корректность аналитических
формул и правильность реализации обоих вариантов нумерации. Отметим также, что
вероятность $P_{\text{both}}$ совпадает для обоих вариантов нумерации (отличие
составляет $4.4 \times 10^{-5}$), что следует из комбинаторного смысла: выбор
двух чётных элементов из множества не зависит от последующей перенумерации.


\clearpage

\section{Моделирование эпидемии на графе социальных контактов}

Требуется реализовать стохастическую модель распространения инфекции
контактным путём на неориентированном графе знакомств. Для каждого контакта
«здоровый–инфицированный» заражение происходит с вероятностью $p_1 \in (0; 1]$.
Инфицированный выздоравливает на каждом шаге с вероятностью $p_2 \in [0; 1]$.
Программа должна предоставлять загрузку графа, пошаговое моделирование с
визуализацией на графе и аналитические запросы к результатам.


\subsection{Модель и структуры данных}

Граф контактов хранится в виде списка смежности. Каждая вершина — структура
\texttt{Person} с полями: идентификатор, состояние (Healthy, Infected,
Recovered) и вектор идентификаторов соседей.


Состояния образуют конечный автомат: Healthy $\to$ Infected (с вероятностью
$p_1$ при контакте), Infected $\to$ Recovered (с вероятностью $p_2$ на каждом
шаге). Состояние Recovered — поглощающее. Повторное заражение невозможно.


\begin{lstlisting}[language=C++, caption={Структура вершины и состояния эпидемиологической модели}, label=lst.person_struct]
enum class PersonState { Healthy, Infected, Recovered };

struct Person {
  int id;
  PersonState state = PersonState::Healthy;
  std::vector<int> contacts;
};
\end{lstlisting}

Генератор псевдослучайных чисел — Вихрь Мерсенна (\texttt{std::mt19937})
с равномерным распределением на $[0, 1)$. На каждом шаге для всех
инфицированных перебираются соседи, и заражение происходит при
$U(0,1) < p_1$. Затем аналогично проверяется выздоровление.


\subsection{Ядро симуляции}

\begin{lstlisting}[language=C++, caption={Метод step() — ядро симуляции}, label=lst.simulation_core]
const std::vector<Person>& Simulation::step() {
    if (m_current_step == 0) {
        std::uniform_int_distribution<size_t> pick(0, m_people.size() - 1);
        m_people[pick(m_rng)].state = PersonState::Infected;
        m_current_step++;
        return m_people;
    }

    std::vector<int> newly_infected;
    for (const auto& person : m_people) {
        if (person.state != PersonState::Infected) continue;
        for (int cid : person.contacts) {
            if ((m_people[cid].state == PersonState::Recovered ||
                 m_people[cid].state == PersonState::Healthy) &&
                m_dist(m_rng) < p_infect)
                newly_infected.push_back(cid);
        }
    }
    for (int id : newly_infected) m_people[id].state = PersonState::Infected;

    for (auto& person : m_people)
        if (person.state == PersonState::Infected && m_dist(m_rng) < p_recover)
            person.state = PersonState::Recovered;

    m_current_step++;
    return m_people;
}
\end{lstlisting}

\subsection{Аналитические запросы к результатам}

Реализованы четыре поисковых запроса, выполняемых за один проход
по массиву вершин:


\begin{itemize}
  \item \textbf{Все не заразившиеся} — вершины, оставшиеся в состоянии Healthy.
  \item \textbf{Все исцелившиеся} — вершины в состоянии Recovered.
  \item \textbf{Исцелившиеся, чьё окружение не исцелилось} — Recovered-вершины, среди соседей которых есть хотя бы одна не-Recovered.
  \item \textbf{Не заразившиеся, чьё окружение полностью заражено} — Healthy-вершины, у которых все соседи имеют состояние Infected.
\end{itemize}

\begin{lstlisting}[language=C++, caption={Поиск исцелившихся с неисцелившимся окружением}, label=lst.query_recovered_lonely]
std::vector<int> Simulation::recovered_with_sick_contacts() const {
    std::vector<int> r;
    for (const auto& p : m_people) {
        if (p.state != PersonState::Recovered) continue;
        for (int cid : p.contacts) {
            if (m_people[cid].state != PersonState::Recovered) {
                r.push_back(p.id);
                break;
            }
        }
    }
    return r;
}
\end{lstlisting}

\subsection{Тестовые данные и визуализация}

Тестовый граф — датасет Facebook-страниц категории «Artist» из коллекции
GEMSEC. Исходный CSV-файл содержит рёбра в формате
\texttt{"u,v"}. Для загрузки выполняется предобработка: запятая заменяется
пробелом. Параметры моделирования: $p_1 = 0.3$, $p_2 = 0.1$.


Визуализация графа реализована на QML с использованием \texttt{Canvas}.
Для работы с большими графами (до $10^5$ вершин) строится BFS-подграф от
вершины с максимальной степенью, ограниченный
заданным числом узлов. Размещение вершин — force-based алгоритм.


Отрисовка дифференцирует состояния вершин цветом: зелёный — Healthy,
красный — Infected, синий — Recovered. Рёбра отображаются при достаточном
масштабе. Поддерживаются панорамирование и масштабирование. Управление
симуляцией и запросами осуществляется через боковую панель интерфейса.


\begin{figure}[h]
\centering
\begin{subfigure}[b]{0.95\linewidth}
\includegraphics[width=\linewidth]{assets/task02_gui.png}
\caption{Интерфейс приложения}
\end{subfigure}
\hfill
\begin{subfigure}[b]{0.95\linewidth}
\includegraphics[width=\linewidth]{assets/task02_search.png}
\caption{Пример аналитического запроса}
\end{subfigure}
\hfill
\caption{Интерфейс приложения и пример аналитического запроса}\label{fig.epidemic_interface}
\end{figure}

\clearpage

\section{Распространение слухов}

В городе проживает $n+1$ человек. Один из них узнаёт новость и передаёт её
другому, тот — следующему, и так далее. На каждом шаге каждый знающий новость
выбирает $N$ случайных адресатов из $n$ остальных жителей и сообщает новость им.
Требуется определить вероятность того, что новость будет передана $r$ раз без:
а) возвращения к первому узнавшему (режим \texttt{ReturnToSender});
б) повторного сообщения кому-либо (режим \texttt{RepeatedMessage}).
Необходимо реализовать приложение для моделирования процесса с настраиваемыми
параметрами $n$, $r$, $N$, $K$ (число экспериментов), поддерживающее пошаговый
и пакетный режимы, визуализацию графа передачи и подсчёт эмпирических вероятностей.


\subsection{Модель и алгоритм}

Процесс моделируется на полном графе из $n+1$ вершин. Вершина $0$ — источник.
Состояние системы — булев массив \texttt{knows} размера $n+1$, изначально
\texttt{knows[0] = true}. На шаге $t$ каждый знающий выбирает $N$ случайных
получателей среди вершин, отличных от него самого.


\subsection{Ядро симуляции}

\begin{lstlisting}[language=C++, caption={Метод run() — выполнение одного эксперимента}, label=lst.rumor_simulation]
std::pair<bool, std::vector<Simulation::GraphNode>> Simulation::run(size_t r, int seed) {
  m_rng = std::mt19937(seed);
  std::uniform_int_distribution<int> dist(0, m_n);
  std::vector<bool> knows(m_n + 1, false);
  knows[0] = true;
  std::vector<GraphNode> graph;

  for (size_t step = 0; step < r; ++step) {
    std::vector<size_t> knowers;
    for (size_t i = 0; i <= m_n; ++i)
      if (knows[i]) knowers.push_back(i);

    for (size_t sender : knowers) {
      for (size_t j = 0; j < m_N; ++j) {
        int target;
        do { target = dist(m_rng); } while (static_cast<size_t>(target) == sender);

        if (m_mode == Mode::ReturnToSender) {
          if (target == 0) {
            graph.push_back({sender, 0, step});
            return {false, graph};
          }
        } else {
          if (knows[static_cast<size_t>(target)]) {
            graph.push_back({sender, static_cast<size_t>(target), step});
            return {false, graph};
          }
        }
        graph.push_back({sender, static_cast<size_t>(target), step});
        knows[static_cast<size_t>(target)] = true;
      }
    }
  }
  return {true, graph};
}
\end{lstlisting}

\subsection{Визуализация и результаты}

Граф передачи для каждого эксперимента отображается в компоненте \\
\texttt{ExperimentGraphView}. Вершины размещаются силовым алгоритмом с анимацией сходимости.


Интерфейс позволяет настраивать $n$, $r$, $N$, $K$, выбирать режим
и запускать моделирование в пошаговом (с визуализацией каждого эксперимента)
или быстром пакетном режиме. Статистика (число успехов/неудач, эмпирическая
вероятность) обновляется в реальном времени.


\begin{figure}[h]
\centering
\begin{subfigure}[b]{0.7\linewidth}
\includegraphics[width=\linewidth]{assets/task03_gui.png}
\caption{Интерфейс приложения: настройка параметров, статистика, граф передачи}
\end{subfigure}
\hfill
\begin{subfigure}[b]{0.7\linewidth}
\includegraphics[width=\linewidth]{assets/task03_success.png}
\caption{Граф успешного эксперимента: новость передана $r$ раз без нарушений}
\end{subfigure}
\hfill
\begin{subfigure}[b]{0.7\linewidth}
\includegraphics[width=\linewidth]{assets/task03_failure.png}
\caption{Граф неудачного эксперимента: нарушение условия (возврат к источнику / повтор)}
\end{subfigure}
\hfill
\caption{Интерфейс приложения и примеры результатов моделирования}\label{fig.rumor_results}
\end{figure}

\clearpage

\section{Моделирование движения материальной точки со случайным шагом}

h


\section{Кластеры и k-паттерны}

Здесь будет описание задачи и её решения.


\section{Ступенчатая фигура}

Здесь будет описание задачи и её решения.


\section{Комплексное блуждание}

Здесь будет описание задачи и её решения.


\section{Бочка с бензином}

Здесь будет описание задачи и её решения.


\section{Криптоанализ}

Здесь будет описание задачи и её решения.


\section{Подбрасывание монеты из кошелька}

Здесь будет описание задачи и её решения.


\section{Дискретная случайная величина}

Здесь будет описание задачи и её решения.


\section{Игла Бюффона}

Здесь будет описание задачи и её решения.


\section{Cтатистика графа}

Здесь будет описание задачи и её решения.


\addcontentsline{toc}{section}{Список использованных источников}

\begin{thebibliography}{1}
  \bibitem{ref.example}
  Фамилия И.О. Название книги. --- Город: Издательство, 2024. 256 с.
\end{thebibliography}


\end{document}
